\documentclass[a4paper,11pt]{article}
\usepackage{operationshandbook}
\begin{document}
\begin{center}
{\Large\bfseries 専門確認・公開承認 業務ハンドブック（独立試案・統合版）}\par\medskip
{\small 独立設計・試作・内容確認未了}\par
{\small v0.2.0\quad 2026年10月19日}\par
\end{center}
\vspace{6mm}
\noindent 本書は組織から承認を受けた規程又は公式の業務マニュアルを示すものではありません。実務への適用に際しては、必要な法令、利用許諾及び組織の権限規程を確認してください。\par
\tableofcontents
\clearpage
\section{目的、適用範囲と確認権限}
\par 本書は、科学的専門確認、編集上の判断、著作物と肖像の権利、個人情報と倫理、公開決裁をそれぞれ別の検討対象として記録する独立の業務試案である。審査対象の版、ファイルハッシュ、媒体、確認担当者、根拠、対象外事項、判断条件を識別する。\par\medskip
\par 実際の権限、契約、適用法令を確かめる必要がある。チェックリストの完了、AIの指摘がないこと、PDFの正常な組版は、適法性・真実性・公開承認の証明にはならない。\par\medskip
\section{審査対象と版の確定}
\subsection{審査対象と版の確定}
\par 担当：審査調整担当者\par\medskip
\subsubsection*{入力資料}
\begin{enumerate}[leftmargin=2em,itemsep=0.35em]
\item 審査対象資料（固定版・ファイルハッシュ）
\item 適用する根拠台帳・素材台帳・確認記録
\item 権限及び公開条件に関する関係記録
\end{enumerate}
\subsubsection*{成果物}
\begin{enumerate}[leftmargin=2em,itemsep=0.35em]
\item 審査対象と版の確定に関する確認記録
\item 未解決事項、差戻し条件及び責任者一覧
\end{enumerate}
\subsubsection*{完了の確認事項}
\begin{enumerate}[leftmargin=2em,itemsep=0.35em]
\item 対象文書の版・ハッシュ・対象媒体・確認者を記録した
\item 担当範囲内の審査結果、証拠、対象外及び未解決事項を識別した
\item 判断保留・差戻し・条件付き判断について責任者と再審査条件を記載した
\end{enumerate}
\par 人による判断を要する：はい\par\medskip
\subsubsection*{実施手順}
\begin{enumerate}[leftmargin=2em,itemsep=0.65em]
\item 対象資料の識別子、版、担当者、受領日時を記録する\par\noindent\textit{記録・成果：}受領記録\nobreak\smallskip
\item 対象版、添付資料及び期限を確定する\par\noindent\textit{記録・成果：}確認結果\nobreak\smallskip
\item 判断根拠、反対意見、未解決事項及び相談先を記録する\par\noindent\textit{記録・成果：}論点記録\nobreak\smallskip
\item 修正の要否、再審査条件、引継ぎ先を明示する\par\noindent\textit{記録・成果：}差戻し・引継ぎ記録\nobreak\smallskip
\end{enumerate}
\subsubsection*{例外・差戻し}
\begin{enumerate}[leftmargin=2em,itemsep=0.35em]
\item 権限者又は根拠が不足する場合は判断を保留する
\item 公開停止を要する疑義は責任者に報告する
\end{enumerate}
\section{審査権限と利益相反の確認}
\subsection{審査権限と利益相反の確認}
\par 担当：審査責任者\par\medskip
\subsubsection*{入力資料}
\begin{enumerate}[leftmargin=2em,itemsep=0.35em]
\item 審査対象資料（固定版・ファイルハッシュ）
\item 適用する根拠台帳・素材台帳・確認記録
\item 権限及び公開条件に関する関係記録
\end{enumerate}
\subsubsection*{成果物}
\begin{enumerate}[leftmargin=2em,itemsep=0.35em]
\item 審査権限と利益相反の確認に関する確認記録
\item 未解決事項、差戻し条件及び責任者一覧
\end{enumerate}
\subsubsection*{完了の確認事項}
\begin{enumerate}[leftmargin=2em,itemsep=0.35em]
\item 対象文書の版・ハッシュ・対象媒体・確認者を記録した
\item 担当範囲内の審査結果、証拠、対象外及び未解決事項を識別した
\item 判断保留・差戻し・条件付き判断について責任者と再審査条件を記載した
\end{enumerate}
\par 人による判断を要する：はい\par\medskip
\subsubsection*{実施手順}
\begin{enumerate}[leftmargin=2em,itemsep=0.65em]
\item 対象資料の識別子、版、担当者、受領日時を記録する\par\noindent\textit{記録・成果：}受領記録\nobreak\smallskip
\item 専門確認、権利確認、倫理確認、公開決裁の権限を分離する\par\noindent\textit{記録・成果：}確認結果\nobreak\smallskip
\item 判断根拠、反対意見、未解決事項及び相談先を記録する\par\noindent\textit{記録・成果：}論点記録\nobreak\smallskip
\item 修正の要否、再審査条件、引継ぎ先を明示する\par\noindent\textit{記録・成果：}差戻し・引継ぎ記録\nobreak\smallskip
\end{enumerate}
\subsubsection*{例外・差戻し}
\begin{enumerate}[leftmargin=2em,itemsep=0.35em]
\item 権限者又は根拠が不足する場合は判断を保留する
\item 公開停止を要する疑義は責任者に報告する
\end{enumerate}
\section{専門的妥当性と主張範囲の審査}
\subsection{専門的妥当性と主張範囲の審査}
\par 担当：科学専門確認者\par\medskip
\subsubsection*{入力資料}
\begin{enumerate}[leftmargin=2em,itemsep=0.35em]
\item 審査対象資料（固定版・ファイルハッシュ）
\item 適用する根拠台帳・素材台帳・確認記録
\item 権限及び公開条件に関する関係記録
\end{enumerate}
\subsubsection*{成果物}
\begin{enumerate}[leftmargin=2em,itemsep=0.35em]
\item 専門的妥当性と主張範囲の審査に関する確認記録
\item 未解決事項、差戻し条件及び責任者一覧
\end{enumerate}
\subsubsection*{完了の確認事項}
\begin{enumerate}[leftmargin=2em,itemsep=0.35em]
\item 対象文書の版・ハッシュ・対象媒体・確認者を記録した
\item 担当範囲内の審査結果、証拠、対象外及び未解決事項を識別した
\item 判断保留・差戻し・条件付き判断について責任者と再審査条件を記載した
\end{enumerate}
\par 人による判断を要する：はい\par\medskip
\subsubsection*{実施手順}
\begin{enumerate}[leftmargin=2em,itemsep=0.65em]
\item 対象資料の識別子、版、担当者、受領日時を記録する\par\noindent\textit{記録・成果：}受領記録\nobreak\smallskip
\item 研究方法、推論、限界及び将来見通しの表現を照合する\par\noindent\textit{記録・成果：}確認結果\nobreak\smallskip
\item 判断根拠、反対意見、未解決事項及び相談先を記録する\par\noindent\textit{記録・成果：}論点記録\nobreak\smallskip
\item 修正の要否、再審査条件、引継ぎ先を明示する\par\noindent\textit{記録・成果：}差戻し・引継ぎ記録\nobreak\smallskip
\end{enumerate}
\subsubsection*{例外・差戻し}
\begin{enumerate}[leftmargin=2em,itemsep=0.35em]
\item 権限者又は根拠が不足する場合は判断を保留する
\item 公開停止を要する疑義は責任者に報告する
\end{enumerate}
\section{命題・根拠・変更版の再照合}
\subsection{命題・根拠・変更版の再照合}
\par 担当：根拠検証担当者\par\medskip
\subsubsection*{入力資料}
\begin{enumerate}[leftmargin=2em,itemsep=0.35em]
\item 審査対象資料（固定版・ファイルハッシュ）
\item 適用する根拠台帳・素材台帳・確認記録
\item 権限及び公開条件に関する関係記録
\end{enumerate}
\subsubsection*{成果物}
\begin{enumerate}[leftmargin=2em,itemsep=0.35em]
\item 命題・根拠・変更版の再照合に関する確認記録
\item 未解決事項、差戻し条件及び責任者一覧
\end{enumerate}
\subsubsection*{完了の確認事項}
\begin{enumerate}[leftmargin=2em,itemsep=0.35em]
\item 対象文書の版・ハッシュ・対象媒体・確認者を記録した
\item 担当範囲内の審査結果、証拠、対象外及び未解決事項を識別した
\item 判断保留・差戻し・条件付き判断について責任者と再審査条件を記載した
\end{enumerate}
\par 人による判断を要する：はい\par\medskip
\subsubsection*{実施手順}
\begin{enumerate}[leftmargin=2em,itemsep=0.65em]
\item 主張の原稿位置と原資料の該当箇所を対応させる\par\noindent\textit{記録・成果：}命題・根拠対応表\nobreak\smallskip
\item 撤回・訂正・新規公開情報の疑義を調査し状態を記録する\par\noindent\textit{記録・成果：}更新確認記録\nobreak\smallskip
\item 修正後の命題と改稿前の命題の差を検証する\par\noindent\textit{記録・成果：}改稿影響表\nobreak\smallskip
\item 図表、短文、字幕、タイトルにも同じ品質の照合を行う\par\noindent\textit{記録・成果：}媒体別確認記録\nobreak\smallskip
\item 再検査未了の命題を未解決として申請対象へ渡す\par\noindent\textit{記録・成果：}未確認命題一覧\nobreak\smallskip
\end{enumerate}
\subsubsection*{例外・差戻し}
\begin{enumerate}[leftmargin=2em,itemsep=0.35em]
\item 原資料へのアクセス不能は検証済みとせず条件付きの照会へ移す
\item 重大な不一致は対象版の再編集へ差し戻す
\end{enumerate}
\section{著作物・肖像・素材利用条件の確認}
\subsection{著作物・肖像・素材利用条件の確認}
\par 担当：権利確認担当者\par\medskip
\subsubsection*{入力資料}
\begin{enumerate}[leftmargin=2em,itemsep=0.35em]
\item 審査対象資料（固定版・ファイルハッシュ）
\item 適用する根拠台帳・素材台帳・確認記録
\item 権限及び公開条件に関する関係記録
\end{enumerate}
\subsubsection*{成果物}
\begin{enumerate}[leftmargin=2em,itemsep=0.35em]
\item 著作物・肖像・素材利用条件の確認に関する確認記録
\item 未解決事項、差戻し条件及び責任者一覧
\end{enumerate}
\subsubsection*{完了の確認事項}
\begin{enumerate}[leftmargin=2em,itemsep=0.35em]
\item 対象文書の版・ハッシュ・対象媒体・確認者を記録した
\item 担当範囲内の審査結果、証拠、対象外及び未解決事項を識別した
\item 判断保留・差戻し・条件付き判断について責任者と再審査条件を記載した
\end{enumerate}
\par 人による判断を要する：はい\par\medskip
\subsubsection*{実施手順}
\begin{enumerate}[leftmargin=2em,itemsep=0.65em]
\item 対象資料の識別子、版、担当者、受領日時を記録する\par\noindent\textit{記録・成果：}受領記録\nobreak\smallskip
\item 利用主体、媒体、期間及び改変条件を資料ごとに確認する\par\noindent\textit{記録・成果：}確認結果\nobreak\smallskip
\item 判断根拠、反対意見、未解決事項及び相談先を記録する\par\noindent\textit{記録・成果：}論点記録\nobreak\smallskip
\item 修正の要否、再審査条件、引継ぎ先を明示する\par\noindent\textit{記録・成果：}差戻し・引継ぎ記録\nobreak\smallskip
\end{enumerate}
\subsubsection*{例外・差戻し}
\begin{enumerate}[leftmargin=2em,itemsep=0.35em]
\item 権限者又は根拠が不足する場合は判断を保留する
\item 公開停止を要する疑義は責任者に報告する
\end{enumerate}
\section{個人情報・機密・公開範囲の確認}
\subsection{個人情報・機密・公開範囲の確認}
\par 担当：情報保護確認担当者\par\medskip
\subsubsection*{入力資料}
\begin{enumerate}[leftmargin=2em,itemsep=0.35em]
\item 審査対象資料（固定版・ファイルハッシュ）
\item 適用する根拠台帳・素材台帳・確認記録
\item 権限及び公開条件に関する関係記録
\end{enumerate}
\subsubsection*{成果物}
\begin{enumerate}[leftmargin=2em,itemsep=0.35em]
\item 個人情報・機密・公開範囲の確認に関する確認記録
\item 未解決事項、差戻し条件及び責任者一覧
\end{enumerate}
\subsubsection*{完了の確認事項}
\begin{enumerate}[leftmargin=2em,itemsep=0.35em]
\item 対象文書の版・ハッシュ・対象媒体・確認者を記録した
\item 担当範囲内の審査結果、証拠、対象外及び未解決事項を識別した
\item 判断保留・差戻し・条件付き判断について責任者と再審査条件を記載した
\end{enumerate}
\par 人による判断を要する：はい\par\medskip
\subsubsection*{実施手順}
\begin{enumerate}[leftmargin=2em,itemsep=0.65em]
\item 対象資料の識別子、版、担当者、受領日時を記録する\par\noindent\textit{記録・成果：}受領記録\nobreak\smallskip
\item 公開目的、情報最小化、アクセス権及び相談先を確認する\par\noindent\textit{記録・成果：}確認結果\nobreak\smallskip
\item 判断根拠、反対意見、未解決事項及び相談先を記録する\par\noindent\textit{記録・成果：}論点記録\nobreak\smallskip
\item 修正の要否、再審査条件、引継ぎ先を明示する\par\noindent\textit{記録・成果：}差戻し・引継ぎ記録\nobreak\smallskip
\end{enumerate}
\subsubsection*{例外・差戻し}
\begin{enumerate}[leftmargin=2em,itemsep=0.35em]
\item 権限者又は根拠が不足する場合は判断を保留する
\item 公開停止を要する疑義は責任者に報告する
\end{enumerate}
\section{倫理・安全・社会的影響の確認}
\subsection{倫理・安全・社会的影響の確認}
\par 担当：倫理確認担当者\par\medskip
\subsubsection*{入力資料}
\begin{enumerate}[leftmargin=2em,itemsep=0.35em]
\item 審査対象資料（固定版・ファイルハッシュ）
\item 適用する根拠台帳・素材台帳・確認記録
\item 権限及び公開条件に関する関係記録
\end{enumerate}
\subsubsection*{成果物}
\begin{enumerate}[leftmargin=2em,itemsep=0.35em]
\item 倫理・安全・社会的影響の確認に関する確認記録
\item 未解決事項、差戻し条件及び責任者一覧
\end{enumerate}
\subsubsection*{完了の確認事項}
\begin{enumerate}[leftmargin=2em,itemsep=0.35em]
\item 対象文書の版・ハッシュ・対象媒体・確認者を記録した
\item 担当範囲内の審査結果、証拠、対象外及び未解決事項を識別した
\item 判断保留・差戻し・条件付き判断について責任者と再審査条件を記載した
\end{enumerate}
\par 人による判断を要する：はい\par\medskip
\subsubsection*{実施手順}
\begin{enumerate}[leftmargin=2em,itemsep=0.65em]
\item 対象資料の識別子、版、担当者、受領日時を記録する\par\noindent\textit{記録・成果：}受領記録\nobreak\smallskip
\item 尊厳、差別、誤認、危害及び当事者への影響を検討する\par\noindent\textit{記録・成果：}確認結果\nobreak\smallskip
\item 判断根拠、反対意見、未解決事項及び相談先を記録する\par\noindent\textit{記録・成果：}論点記録\nobreak\smallskip
\item 修正の要否、再審査条件、引継ぎ先を明示する\par\noindent\textit{記録・成果：}差戻し・引継ぎ記録\nobreak\smallskip
\end{enumerate}
\subsubsection*{例外・差戻し}
\begin{enumerate}[leftmargin=2em,itemsep=0.35em]
\item 権限者又は根拠が不足する場合は判断を保留する
\item 公開停止を要する疑義は責任者に報告する
\end{enumerate}
\section{専門指摘・修正・異議の調整と再審}
\subsection{専門指摘・修正・異議の調整と再審}
\par 担当：審査調整担当者\par\medskip
\subsubsection*{入力資料}
\begin{enumerate}[leftmargin=2em,itemsep=0.35em]
\item 審査対象資料（固定版・ファイルハッシュ）
\item 適用する根拠台帳・素材台帳・確認記録
\item 権限及び公開条件に関する関係記録
\end{enumerate}
\subsubsection*{成果物}
\begin{enumerate}[leftmargin=2em,itemsep=0.35em]
\item 専門指摘・修正・異議の調整と再審に関する確認記録
\item 未解決事項、差戻し条件及び責任者一覧
\end{enumerate}
\subsubsection*{完了の確認事項}
\begin{enumerate}[leftmargin=2em,itemsep=0.35em]
\item 対象文書の版・ハッシュ・対象媒体・確認者を記録した
\item 担当範囲内の審査結果、証拠、対象外及び未解決事項を識別した
\item 判断保留・差戻し・条件付き判断について責任者と再審査条件を記載した
\end{enumerate}
\par 人による判断を要する：はい\par\medskip
\subsubsection*{実施手順}
\begin{enumerate}[leftmargin=2em,itemsep=0.65em]
\item 指摘を命題、素材、媒体、重大度別に登録する\par\noindent\textit{記録・成果：}指摘管理台帳\nobreak\smallskip
\item 採用、追加確認、却下、保留の判断理由と責任者を記録する\par\noindent\textit{記録・成果：}対応判断表\nobreak\smallskip
\item 改稿・素材差替えが他媒体や権利に及ぼす影響を検討する\par\noindent\textit{記録・成果：}影響分析表\nobreak\smallskip
\item 重要な修正を対応する専門家・権限者へ再確認する\par\noindent\textit{記録・成果：}再審結果\nobreak\smallskip
\item 未解決の高重大度事項を公開判断のゲートに残す\par\noindent\textit{記録・成果：}未解決事項表\nobreak\smallskip
\end{enumerate}
\subsubsection*{例外・差戻し}
\begin{enumerate}[leftmargin=2em,itemsep=0.35em]
\item 異議を出した専門家自身の意見を編集上の承認として自動確定しない
\item 未解決重大事項は保留し必要な以前の工程に差し戻す
\end{enumerate}
\section{指摘対応と再審査}
\subsection{指摘対応と再審査}
\par 担当：編集責任者\par\medskip
\subsubsection*{入力資料}
\begin{enumerate}[leftmargin=2em,itemsep=0.35em]
\item 審査対象資料（固定版・ファイルハッシュ）
\item 適用する根拠台帳・素材台帳・確認記録
\item 権限及び公開条件に関する関係記録
\end{enumerate}
\subsubsection*{成果物}
\begin{enumerate}[leftmargin=2em,itemsep=0.35em]
\item 指摘対応と再審査に関する確認記録
\item 未解決事項、差戻し条件及び責任者一覧
\end{enumerate}
\subsubsection*{完了の確認事項}
\begin{enumerate}[leftmargin=2em,itemsep=0.35em]
\item 対象文書の版・ハッシュ・対象媒体・確認者を記録した
\item 担当範囲内の審査結果、証拠、対象外及び未解決事項を識別した
\item 判断保留・差戻し・条件付き判断について責任者と再審査条件を記載した
\end{enumerate}
\par 人による判断を要する：はい\par\medskip
\subsubsection*{実施手順}
\begin{enumerate}[leftmargin=2em,itemsep=0.65em]
\item 対象資料の識別子、版、担当者、受領日時を記録する\par\noindent\textit{記録・成果：}受領記録\nobreak\smallskip
\item 重大度、担当者、対応結果及び再確認の要否を明示する\par\noindent\textit{記録・成果：}確認結果\nobreak\smallskip
\item 判断根拠、反対意見、未解決事項及び相談先を記録する\par\noindent\textit{記録・成果：}論点記録\nobreak\smallskip
\item 修正の要否、再審査条件、引継ぎ先を明示する\par\noindent\textit{記録・成果：}差戻し・引継ぎ記録\nobreak\smallskip
\end{enumerate}
\subsubsection*{例外・差戻し}
\begin{enumerate}[leftmargin=2em,itemsep=0.35em]
\item 権限者又は根拠が不足する場合は判断を保留する
\item 公開停止を要する疑義は責任者に報告する
\end{enumerate}
\section{公表日時・共同発信と公開条件の検証}
\subsection{公表日時・共同発信と公開条件の検証}
\par 担当：公開調整担当者\par\medskip
\subsubsection*{入力資料}
\begin{enumerate}[leftmargin=2em,itemsep=0.35em]
\item 審査対象資料（固定版・ファイルハッシュ）
\item 適用する根拠台帳・素材台帳・確認記録
\item 権限及び公開条件に関する関係記録
\end{enumerate}
\subsubsection*{成果物}
\begin{enumerate}[leftmargin=2em,itemsep=0.35em]
\item 公表日時・共同発信と公開条件の検証に関する確認記録
\item 未解決事項、差戻し条件及び責任者一覧
\end{enumerate}
\subsubsection*{完了の確認事項}
\begin{enumerate}[leftmargin=2em,itemsep=0.35em]
\item 対象文書の版・ハッシュ・対象媒体・確認者を記録した
\item 担当範囲内の審査結果、証拠、対象外及び未解決事項を識別した
\item 判断保留・差戻し・条件付き判断について責任者と再審査条件を記載した
\end{enumerate}
\par 人による判断を要する：はい\par\medskip
\subsubsection*{実施手順}
\begin{enumerate}[leftmargin=2em,itemsep=0.65em]
\item 論文・会見・ウェブ等の公開時点とタイムゾーンを照合する\par\noindent\textit{記録・成果：}日時照合記録\nobreak\smallskip
\item 禁止事項の根拠を契約、運用ルール、合意ごとに区別する\par\noindent\textit{記録・成果：}条件根拠表\nobreak\smallskip
\item 共同発表者・関係窓口・公開順序を確認する\par\noindent\textit{記録・成果：}連携確認記録\nobreak\smallskip
\item 公開直前の時刻変更・中止連絡先を定める\par\noindent\textit{記録・成果：}緊急連絡経路\nobreak\smallskip
\item 各媒体・配信システムのスケジュール状態を点検する\par\noindent\textit{記録・成果：}配信条件記録\nobreak\smallskip
\end{enumerate}
\subsubsection*{例外・差戻し}
\begin{enumerate}[leftmargin=2em,itemsep=0.35em]
\item 日時・禁輸の条件が不明確なら配信を保留する
\item 時間帯の変換結果が一致しなければ再確認する
\end{enumerate}
\section{公開決裁と条件付き判断}
\subsection{公開決裁と条件付き判断}
\par 担当：権限ある公開決裁者\par\medskip
\subsubsection*{入力資料}
\begin{enumerate}[leftmargin=2em,itemsep=0.35em]
\item 審査対象資料（固定版・ファイルハッシュ）
\item 適用する根拠台帳・素材台帳・確認記録
\item 権限及び公開条件に関する関係記録
\end{enumerate}
\subsubsection*{成果物}
\begin{enumerate}[leftmargin=2em,itemsep=0.35em]
\item 公開決裁と条件付き判断に関する確認記録
\item 未解決事項、差戻し条件及び責任者一覧
\end{enumerate}
\subsubsection*{完了の確認事項}
\begin{enumerate}[leftmargin=2em,itemsep=0.35em]
\item 対象文書の版・ハッシュ・対象媒体・確認者を記録した
\item 担当範囲内の審査結果、証拠、対象外及び未解決事項を識別した
\item 判断保留・差戻し・条件付き判断について責任者と再審査条件を記載した
\end{enumerate}
\par 人による判断を要する：はい\par\medskip
\subsubsection*{実施手順}
\begin{enumerate}[leftmargin=2em,itemsep=0.65em]
\item 対象資料の識別子、版、担当者、受領日時を記録する\par\noindent\textit{記録・成果：}受領記録\nobreak\smallskip
\item 権限者が承認、差戻し、保留又は却下を理由付きで記録する\par\noindent\textit{記録・成果：}確認結果\nobreak\smallskip
\item 判断根拠、反対意見、未解決事項及び相談先を記録する\par\noindent\textit{記録・成果：}論点記録\nobreak\smallskip
\item 修正の要否、再審査条件、引継ぎ先を明示する\par\noindent\textit{記録・成果：}差戻し・引継ぎ記録\nobreak\smallskip
\end{enumerate}
\subsubsection*{例外・差戻し}
\begin{enumerate}[leftmargin=2em,itemsep=0.35em]
\item 権限者又は根拠が不足する場合は判断を保留する
\item 公開停止を要する疑義は責任者に報告する
\end{enumerate}
\section{責任ある公開可否の判断と版の承認}
\subsection{責任ある公開可否の判断と版の承認}
\par 担当：権限ある公開決裁者\par\medskip
\subsubsection*{入力資料}
\begin{enumerate}[leftmargin=2em,itemsep=0.35em]
\item 審査対象資料（固定版・ファイルハッシュ）
\item 適用する根拠台帳・素材台帳・確認記録
\item 権限及び公開条件に関する関係記録
\end{enumerate}
\subsubsection*{成果物}
\begin{enumerate}[leftmargin=2em,itemsep=0.35em]
\item 責任ある公開可否の判断と版の承認に関する確認記録
\item 未解決事項、差戻し条件及び責任者一覧
\end{enumerate}
\subsubsection*{完了の確認事項}
\begin{enumerate}[leftmargin=2em,itemsep=0.35em]
\item 対象文書の版・ハッシュ・対象媒体・確認者を記録した
\item 担当範囲内の審査結果、証拠、対象外及び未解決事項を識別した
\item 判断保留・差戻し・条件付き判断について責任者と再審査条件を記載した
\end{enumerate}
\par 人による判断を要する：はい\par\medskip
\subsubsection*{実施手順}
\begin{enumerate}[leftmargin=2em,itemsep=0.65em]
\item 判断者の権限と対象媒体・案件における委任範囲を確認する\par\noindent\textit{記録・成果：}権限根拠記録\nobreak\smallskip
\item 科学・編集・権利・個人情報・日時の個別審査状態を確認する\par\noindent\textit{記録・成果：}審査結果一覧\nobreak\smallskip
\item 重大な未解決問題と不適合を分類し対応可能性を判断する\par\noindent\textit{記録・成果：}未解決問題評価\nobreak\smallskip
\item 対象の版・添付素材・媒体範囲と採否理由を明示して決裁する\par\noindent\textit{記録・成果：}決裁記録\nobreak\smallskip
\item 条件・失効・重要変更時の再決裁・中止条件を定める\par\noindent\textit{記録・成果：}承認条件表\nobreak\smallskip
\item 判断の結果を編集・配信担当者へ記録付きで通知する\par\noindent\textit{記録・成果：}承認通知・引継ぎ\nobreak\smallskip
\end{enumerate}
\subsubsection*{例外・差戻し}
\begin{enumerate}[leftmargin=2em,itemsep=0.35em]
\item 審査票の全項目「適合」は自動的な公開承認を意味しない
\item 権限が未確認・重大な権利疑義が残る場合は公開に進めない
\item 条件付きの決裁状態を配信許可へ自動変換しない
\end{enumerate}
\section{公開用パッケージの引継ぎ}
\subsection{公開用パッケージの引継ぎ}
\par 担当：公開工程担当者\par\medskip
\subsubsection*{入力資料}
\begin{enumerate}[leftmargin=2em,itemsep=0.35em]
\item 審査対象資料（固定版・ファイルハッシュ）
\item 適用する根拠台帳・素材台帳・確認記録
\item 権限及び公開条件に関する関係記録
\end{enumerate}
\subsubsection*{成果物}
\begin{enumerate}[leftmargin=2em,itemsep=0.35em]
\item 公開用パッケージの引継ぎに関する確認記録
\item 未解決事項、差戻し条件及び責任者一覧
\end{enumerate}
\subsubsection*{完了の確認事項}
\begin{enumerate}[leftmargin=2em,itemsep=0.35em]
\item 対象文書の版・ハッシュ・対象媒体・確認者を記録した
\item 担当範囲内の審査結果、証拠、対象外及び未解決事項を識別した
\item 判断保留・差戻し・条件付き判断について責任者と再審査条件を記載した
\end{enumerate}
\par 人による判断を要する：はい\par\medskip
\subsubsection*{実施手順}
\begin{enumerate}[leftmargin=2em,itemsep=0.65em]
\item 対象資料の識別子、版、担当者、受領日時を記録する\par\noindent\textit{記録・成果：}受領記録\nobreak\smallskip
\item 承認済み版と条件を配信担当へ渡し再承認条件を記録する\par\noindent\textit{記録・成果：}確認結果\nobreak\smallskip
\item 判断根拠、反対意見、未解決事項及び相談先を記録する\par\noindent\textit{記録・成果：}論点記録\nobreak\smallskip
\item 修正の要否、再審査条件、引継ぎ先を明示する\par\noindent\textit{記録・成果：}差戻し・引継ぎ記録\nobreak\smallskip
\end{enumerate}
\subsubsection*{例外・差戻し}
\begin{enumerate}[leftmargin=2em,itemsep=0.35em]
\item 権限者又は根拠が不足する場合は判断を保留する
\item 公開停止を要する疑義は責任者に報告する
\end{enumerate}
\section{関連する品質標準・運用原則}
\par 以下の標準は独立した試案である。外部のガイドラインや公的解説を参照している場合も、権限を有する規程として採用されたことは意味しない。\par\medskip
\subsection{専門確認の範囲・専門性・独立性}
\subsubsection{審査対象と依頼目的}
\par 専門確認の対象を科学的正確性、方法・推論、専門用語、適用限界、媒体別表現に分けて明示する。対象外となる権利・公開決裁は別の判断へ送る。\par\medskip
\subsubsection{適格性と利益相反}
\par 確認者の分野・経験、関与の範囲、対象研究との関係及び利益相反を記録する。利害関係が存在する場合の独立確認の要否を検討する。\par\medskip
\subsubsection{回答の範囲}
\par 専門家への依頼には対象版、命題ID、回答期限、所望の確認粒度を付す。部分的な確認を全文承認として扱わない。\par\medskip
\subsubsection{異議の扱い}
\par 異なる専門的見解は原資料・根拠とともに保持し、公開判断者へ未解決事項として渡す。\par\medskip
\subsection{確認根拠と対象版の完全性}
\subsubsection{照合可能な版}
\par 本文、画像、動画、字幕、ウェブ断片の識別子・版・ハッシュを固定して確認対象を明示する。\par\medskip
\subsubsection{記録粒度}
\par 指摘にページ・節・時刻・命題IDと原資料への参照を結び付ける。\par\medskip
\subsubsection{更新・撤回}
\par 先行結果の訂正、撤回、更新の疑義について確認日時と調査範囲を記録し、結論を断定し過ぎない。\par\medskip
\subsubsection{再確認}
\par 重要な修正の後は影響した命題、図表、表題について必要な範囲で再確認を依頼する。\par\medskip
\subsection{科学的妥当性・推論範囲の確認}
\subsubsection{命題別点検}
\par 対象記事の命題を研究で測定した事実、統計的関連、因果解釈、外挿、将来の構想に分けて照合する。\par\medskip
\subsubsection{方法と不確実性}
\par 研究設計、母集団、標本、対照、効果量、研究段階、再現性に関する限界を必要な精度で記録する。\par\medskip
\subsubsection{媒体間の伝播}
\par 見出し、図表、字幕、SNS用要約等の強い表現にも同じ主張IDを関連付ける。\par\medskip
\subsubsection{見解の帰属}
\par 専門家が確認した範囲と編集責任者が選択した表現を区別し、根拠のない保証や推奨を付与しない。\par\medskip
\subsection{著作物・画像・音声等の権利調査}
\subsubsection{利用素材台帳}
\par 著作者・権利管理者、利用媒体、改変の有無、引用等の法的根拠、許諾の範囲と期限を素材単位で記録する。\par\medskip
\subsubsection{異なる権利}
\par 著作権、著作者人格権、肖像、名誉、商標、契約、個人情報を一括の許可として扱わない。\par\medskip
\subsubsection{公開方法別の確認}
\par ウェブ、PDF、SNS、動画、二次配布、再編集等の利用態様に応じて既存許諾の適用範囲を確かめる。\par\medskip
\subsubsection{判断と相談}
\par 例外規定の該当性等を断定しにくい案件は権限ある担当又は専門家に照会し、公開方法の変更も検討する。\par\medskip
\subsection{個人情報・機微情報の公開前検査}
\subsubsection{情報の性格}
\par 氏名・肖像・音声・位置・健康関連の情報等について識別可能性と取得・利用状況を個別に把握する。\par\medskip
\subsubsection{適用主体}
\par 対象業務の主体、利用目的、第三者提供その他の具体的事実に基づいて適用法令とガイドラインを調べる。\par\medskip
\subsubsection{必要性と最小化}
\par 公開目的と必要性を照らし、伏字、匿名化、編集、非公開への切替え等の選択肢を比較する。\par\medskip
\subsubsection{再識別と拡散}
\par 単一素材だけでなく字幕、メタデータ、他の公開資料との照合による識別可能性を確認する。\par\medskip
\subsection{公表日時・禁輸・共同発信の調整}
\subsubsection{時点の特定}
\par 論文公開、会見、配信、各媒体の時刻を時区分付きで記録する。日本時間への換算や日付変更に注意する。\par\medskip
\subsubsection{根拠と制約}
\par 公開制限が契約、出版社の運用、当事者の合意など何に基づくかを区別し、期限切れ情報を更新する。\par\medskip
\subsubsection{共同発信}
\par 複数の発表主体や制作会社がある場合は版・広報窓口・公開チャネル・更新責任を合意した内容に従い確認する。\par\medskip
\subsubsection{変更時の扱い}
\par 差止め、時刻変更、誤配信を想定し、連絡経路・撤回手段・必要な判断者を決める。\par\medskip
\subsection{公開判断と承認権限の分離}
\subsubsection{判断の種類}
\par 科学的確認、法令・権利判断、倫理評価、編集評価、公開決裁は目的と責任が異なる。判断を統合して単一の無条件合格とみなさない。\par\medskip
\subsubsection{権限確認}
\par 公開決裁を行う人物または機関と決裁範囲を実際の規程・委任に照らして特定する。モデル上の役割名だけで権限は生じない。\par\medskip
\subsubsection{状態遷移}
\par 承認、条件付き保留、差戻し、不承認、追加照会の区分を設け、条件付き承認の自動公開を禁止する。\par\medskip
\subsubsection{記録と再審}
\par 決裁対象版、承認者、日付、理由、付帯条件、再審要件を保存し、重要な変更は必要な承認を再取得する。\par\medskip
\subsection{公開用パッケージ・最終受渡し条件}
\subsubsection{公開版の固定}
\par 記事・添付素材・字幕・説明文・掲載先の組合せを列挙し、承認対象の版と送信版の同一性を確認する。\par\medskip
\subsubsection{媒体別準備}
\par アクセシビリティ、問い合わせ先、出典、訂正経路、コンテンツ警告等を対象媒体に応じて確認する。\par\medskip
\subsubsection{差戻し条件}
\par 未解消の権利・人権・科学的誤認や公開時刻不一致がある場合は公開に移さず判断者へ返す。\par\medskip
\subsubsection{引継ぎ記録}
\par 公開担当者へ配信可能版、公開日時、承認記録、撤回条件、監視窓口を渡し受領を記録する。\par\medskip
\section{倫理的判断と関係主体への配慮}
\subsection{専門審査の独立性と利益相反}
\subsubsection{判断の尊重}
\par 研究者・制作担当・編集責任者の評価範囲を明示し、説明内容が検証対象にあることを尊重する。\par\medskip
\subsubsection{利益相反の扱い}
\par 利害関係、役務受託、研究資金等の関係を必要に応じて開示し、審査体制の独立性を検討する。\par\medskip
\subsubsection{異論の保存}
\par 異論・少数意見を根拠付きで記録し、都合のよい部分だけを「専門家確認済み」と表示しない。\par\medskip
\subsection{被取材者の尊厳・私生活と同意の境界}
\subsubsection{本人の利益}
\par 発言、画像、取材記録の公開が本人や周囲へもたらす影響を具体的に想定する。\par\medskip
\subsubsection{同意の区別}
\par 取材参加、録音、引用、画像利用、公開対象・目的の説明は別の事項として評価する。同意のみをすべての権利・法的根拠とみなさない。\par\medskip
\subsubsection{適切な相談}
\par センシティブな情報、未成年者、再識別のおそれがある場合は必要な専門的確認と公開範囲の再検討を行う。\par\medskip
\subsection{予見可能な社会的被害と表現上の配慮}
\subsubsection{有害な読み違い}
\par 健康、教育、社会的属性等のテーマでは将来予測や効果の誇張が不当な期待や危険を生む可能性を検討する。\par\medskip
\subsubsection{文脈保持}
\par 研究の限界、対象集団、技術水準と代替的解釈を、受け手の理解に必要な範囲で提供する。\par\medskip
\subsubsection{修正と救済}
\par 不利益の指摘を受ける窓口を整え、訂正・非公開化・追加説明の判断経路を記録する。\par\medskip
\subsection{AI支援・外部提供と人による責任}
\subsubsection{データの入力範囲}
\par 未公開原稿、個人情報、画像、契約、評価コメントを外部AIへ送る前に契約・内部基準・アクセス条件を確認する。\par\medskip
\subsubsection{機械判定の限界}
\par AIが検出した疑義や推奨を人の確認候補として扱い、専門確認や承認の代行とみなさない。\par\medskip
\subsubsection{記録}
\par AI利用の範囲、入力の類型、生成結果、確認者による採否と修正根拠を記録する。\par\medskip
\section{専門確認チェックリスト}
\subsection{専門確認チェックリスト}
\par 確認用の定義であり、個別案件での実施状況、根拠資料、確認者、レビューの結果は別途保存する。\par\medskip
\subsubsection*{審査範囲・根拠}
\begin{enumerate}[leftmargin=2em,itemsep=0.3em]
\item 確認対象の版、ハッシュ、媒体と審査範囲が特定されている\par
\item 確認者の専門性、担当範囲と利益相反が記録されている\par
\item 主要な科学的主張が原資料及び研究方法に遡れる\par
\item 因果解釈、一般化、数値、限界の表現を点検した\par
\end{enumerate}
\subsubsection*{対応判断・記録}
\begin{enumerate}[leftmargin=2em,itemsep=0.3em]
\item 未公表・撤回・訂正情報の影響を確認した\par
\item 対象外の専門領域と未確認の問題を示した\par
\item 指摘と対応担当者、対応期限が紐付いている\par
\item 重大な修正は専門家による再確認対象にした\par
\end{enumerate}
\section{レビュー指摘対応チェックリスト}
\subsection{レビュー指摘対応チェックリスト}
\par 確認用の定義であり、個別案件での実施状況、根拠資料、確認者、レビューの結果は別途保存する。\par\medskip
\subsubsection*{審査範囲・根拠}
\begin{enumerate}[leftmargin=2em,itemsep=0.3em]
\item 全ての指摘に固有ID、発信者、対象版がある\par
\item 重要度、科学的影響、公開上の影響を区別した\par
\item 変更前後の差分と判断理由が参照可能である\par
\item 採用しなかった指摘と異論を保存した\par
\end{enumerate}
\subsubsection*{対応判断・記録}
\begin{enumerate}[leftmargin=2em,itemsep=0.3em]
\item 修正により他の媒体や図表へ影響が出たか確認した\par
\item 再審査が必要な項目と再確認者を特定した\par
\item 未完了指摘が公開決裁者へ明示されている\par
\item 指摘の閉鎖が権限ある担当者により記録された\par
\end{enumerate}
\section{公開承認ゲートチェックリスト}
\subsection{公開承認ゲートチェックリスト}
\par 確認用の定義であり、個別案件での実施状況、根拠資料、確認者、レビューの結果は別途保存する。\par\medskip
\subsubsection*{審査範囲・根拠}
\begin{enumerate}[leftmargin=2em,itemsep=0.3em]
\item 全媒体の公開対象版及びハッシュが固定されている\par
\item 専門、編集、権利、情報保護、倫理の確認が個別に記録された\par
\item 各確認の未解決事項、対象外事項及び留保が公開決裁者に提示された\par
\item 公表日時、禁輸、共同発信の条件が参照可能である\par
\item 承認者の職務権限と当該案件への適用を別途確認した\par
\end{enumerate}
\subsubsection*{対応判断・記録}
\begin{enumerate}[leftmargin=2em,itemsep=0.3em]
\item 判断区分、理由、条件、対象版、日時が一つの記録に紐付く\par
\item 条件付き承認を無条件承認として扱っていない\par
\item 公開直前の変更に対する再審査・停止手続がある\par
\item 配信担当者に確定版と公開条件を安全に引き継いだ\par
\item 審査記録やチェック完了を独立した公開許可として扱っていない\par
\end{enumerate}
\section{公開前権利確認チェックリスト}
\subsection{公開前権利確認チェックリスト}
\par 確認用の定義であり、個別案件での実施状況、根拠資料、確認者、レビューの結果は別途保存する。\par\medskip
\subsubsection*{審査範囲・根拠}
\begin{enumerate}[leftmargin=2em,itemsep=0.3em]
\item 画像、動画、音声、図表、文章を素材単位で識別した\par
\item 著作者・権利保有者・作成過程を確認した\par
\item 許諾又は法定例外を当該利用方法ごとに検討した\par
\item 引用の場合は公表状況、必要性、範囲、出所を検討した\par
\end{enumerate}
\subsubsection*{対応判断・記録}
\begin{enumerate}[leftmargin=2em,itemsep=0.3em]
\item 著作者人格権、肖像、商標、個人情報等を別個に検討した\par
\item 編集・翻訳・二次配布・SNS再掲の範囲を確認した\par
\item 期限切れ又は未確認の素材に差替え・保留方針を定めた\par
\item 確認結果と公開方法を権限ある担当者が確認した\par
\end{enumerate}
\section{公開判断、構造検査と証跡の境界}
\par 公開判断は、対象ファイルの内容ハッシュと決裁記録を結び付け、権限と責任のある人が実施する。条件付き、差戻し、保留、却下を無条件承認と同じ状態へ機械的に統合しない。\par\medskip
\par 構造検査プログラムは、必須領域の確認記録、未解決事項、版及び判断記録の整合を検査する。決裁者の実在、署名、授権範囲、法令適合性、公開権限は検証できず、自動公開を許可しない。\par\medskip
\par 科学的正確性には研究者・専門家による審査が、権利及び個人情報には該当する権利や適用条件の個別判断が、公開には権限ある決裁者の判断が必要である。配信後も訂正、取下げ及び再審査に備え記録を保持する。\par\medskip
\section{主要参照資料と出典の性格}
\par 学術的な研究成果の表現：日本心理学会「科学コミュニケーション・ガイドライン」 \url{https://psych.or.jp/jpamember/press_release_guidelines/} 。研究成果発表に関する学会の専門的推奨である。\par\medskip
\par 編集倫理：日本新聞協会「新聞倫理綱領」 \url{https://www.pressnet.or.jp/outline/ethics/} 。独立した報道の倫理を示す職業的参照枠組みであり、組織広報の手続規程と同一ではない。\par\medskip
\par 個人情報：個人情報保護委員会「ガイドライン（通則編）」 \url{https://www.ppc.go.jp/personalinfo/legal/guidelines_tsusoku/} 。対象主体、適用条文、例外と現在の施行内容を案件ごとに確認する。\par\medskip
\par 著作物の利用：文化庁「他人の著作物を利用したい場合など」 \url{https://www.bunka.go.jp/seisaku/chosakuken/seidokaisetsu/chosakukensha_fumei/} 。著作物性、保護期間、権利制限、許諾等を区別して検討する参考資料である。\par\medskip
\end{document}
